\section{Discussion}\label{sec:discussion}

The central empirical pattern is that ordinary revisions of university IP policies are followed by different licensing
outcomes depending on the direction of the revision, while invention disclosure does not move in the same way. In the
preferred documentary-reviewed sample, revisions toward more supportive communication are followed by more licenses and
options relative to clean controls, whereas revisions toward more restrictive communication are followed by fewer. The
average upward-minus-downward post-revision difference is about one-half log point. The same contrast is visible under the
stricter post-support rule and remains positive after dropping any single event.

The expanded sample changes how this result should be read. The September-25 draft relied on 25 hand-reviewed events,
including only six upward revisions. The revised workflow begins from all 127 detected revisions, applies outcome-support
and non-overlap restrictions mechanically, and then subjects every one of the 62 mechanically eligible pairs to the same
documentary comparability and timing rules. This yields 34 events, eight upward and 26 downward. The preferred estimate is
smaller than in the earlier 25-event specification, rather than larger, but remains positive and statistically unusual in
the direction-label permutation distribution. That attenuation is reassuring in the limited sense that the expanded
documentary review does not mechanically select events that strengthen the headline result.

The evidence is nevertheless outcome-specific. Filing and patent applications move in the same direction as licensing but
with substantially less precision, and invention disclosures do not increase following upward revisions. The
Benjamini--Hochberg adjustment across the five main outcomes raises the licensing $q$-value to 0.159. The paper therefore
does not establish that supportive revisions improve technology transfer generally. Instead, the strongest evidence is
that post-revision divergence is concentrated at the licensing margin.

That location fits the literature that treats the TTO as an organization rather than a passive conduit. Licensing is the
stage at which office staffing, routines, market search, negotiation, and relationships with firms matter most
\citep{SiegelWaldmanLink2003,ThursbyKemp2002,MarkmanEtAl2005}. It is also the stage at which continued cooperation among
inventors, TTO staff, and prospective licensees is especially important for embryonic university technologies
\citep{JensenThursby2001}. By contrast, invention disclosure is closer to the faculty-entry margin emphasized by
\citet{JensenThursbyThursby2003}, \citet{OwenSmithPowell2001}, and \citet{BercovitzFeldman2008}. The absence of a disclosure
response is therefore informative about where the empirical pattern appears, even if the data cannot identify the mechanism.

The content evidence is consistent with this organizational interpretation but should remain secondary. Across the full
127-revision universe, downward revisions are more likely to add conflict-of-interest language. Such provisions govern
faculty relationships with firms and can plausibly affect the complexity of licensing negotiations, but keyword presence
does not establish legal meaning or administrative burden. The provision-level coding of the original 25-event subset
similarly suggests differences in release, waiver, appeal, and revenue-sharing provisions, but those labels were not
retroactively extended after observing the expanded-sample results. A future paper with independent provision coding could
test these channels directly.

The timing evidence also bears on the possibility that written policies are ceremonial. Institutional theory allows formal
rules to be adopted for legitimacy and decoupled from practice \citep{MeyerRowan1977}. The observed divergence after the
documented revision dates is difficult to reconcile with complete decoupling, but it does not imply that wording itself
causes behavior. A revision may be the visible record of a broader organizational change: new TTO leadership, altered
commercialization priorities, reorganized workflows, or changes in university-industry relationships. The administrative
burden literature provides one mechanism through which formal procedures could matter
\citep{MoynihanHerdHarvey2015,HerdMoynihan2018}, but the present design cannot separate textual effects from organizational
changes that travel with a rewrite.

The most important identification concern is selection into revision direction. Upward and downward events have similar
pre-revision licensing trends, and the preferred estimate survives the stricter post-support sample and leave-one-out
analysis. Yet upward revisions begin from markedly less supportive prior documents: prior-document PCSI averages 0.384 for
upward events and 0.485 for downward events. This imbalance may reflect mean reversion, institutional heterogeneity, or
different reasons for revising. Direction-label permutation inference is useful with only eight upward events, but the
exchangeability assumption is substantive because direction was not randomly assigned. Neither flat pre-trends nor
permutation inference rules out contemporaneous direction-specific organizational changes.

For universities and state systems, the policy implication is therefore restrained. Ordinary IP-policy revisions should
not be treated as administrative housekeeping: they coincide with economically meaningful changes in licensing outcomes,
and the direction of communication contains information about where those changes occur. But the estimates do not support
a claim that making policy language more supportive will mechanically raise licensing by a predictable amount. A stronger
causal design would exploit reforms imposed by a governing system, staggered adoption across campuses, legal shocks with
differential exposure, or other settings in which revision direction is plausibly exogenous.

These conclusions are also bounded by the data. The preferred sample contains 34 events and only eight upward revisions;
some timing classifications remain tier B; AUTM reporting is voluntary; outcomes are institution-level aggregates; and the
survey does not follow individual inventions from disclosure through licensing. The 2023 endpoint also limits post-revision
support for the latest events, which is why the 29-event stricter sample is reported prominently. Within these limits, the
evidence is consistent and specific: post-revision divergence is strongest at licensing, not at disclosure.

\section{Conclusion}\label{sec:conclusion}

This paper asks what follows when universities rewrite their intellectual-property policies. Starting from 127 detected
changes in observed policy histories, it applies common mechanical eligibility rules and document-level review to build a
preferred sample of 34 dated, comparable revisions. In a stacked event study, revisions toward more supportive
communication are followed by more licensing than revisions toward more restrictive communication, with no statistically
detectable differential pre-trend. The preferred licensing gap is 0.495 log points; a stricter 29-event sample gives 0.499.
Filing and patent applications move in the same direction but imprecisely, while invention disclosures do not increase.

The contribution is not a causal estimate of wording. It is evidence on the ordinary within-university revisions that
institutions and state systems actually make, evidence on where along the technology-transfer pipeline post-revision change
concentrates, and a transparent method for turning incomplete policy archives into auditable event-study samples. The
expanded documentary review strengthens the empirical foundation while also making the limits clearer: upward events remain
few, starting policies differ sharply by direction, and universities choose how to revise.

Future work should seek settings in which revision timing or direction is externally induced and should code the legal and
administrative content of revisions independently of the communication measure. Dated, versioned policy archives linked to
campus-level implementation and invention-level outcomes would make it possible to distinguish the effect of policy text
from the organizational changes that accompany it.

\section*{Declarations}
\noindent\textbf{Funding:} \\
\textbf{Conflicts of interest:} The authors declare no conflicts of interest.\\
\textbf{Data availability:} The policy corpus, PCSI scores, revision codes, and replication code are provided in the
replication package. AUTM survey data are provided subject to AUTM's terms of use; users may need to obtain the survey data
directly from AUTM.\\
\textbf{Code availability:} The replication package reproduces the reported sample construction and empirical estimates;
see its \texttt{README}.\\
\textbf{Author contributions:}
